\chapter{Wariacja kwadratowa i kowariacja ciągłych martyngałów}
\label{ch:quadratic-cont}
\paragraph{Odsyłacze międzytomowe.} Konstrukcja całki Itô wykorzystująca wariację kwadratową znajduje się w tomie III, rozdziale~\ref{III-ch:ito-construction}.


W całym rozdziale zakładamy warunki zwykłe filtracji
(rozdział~\ref{ch:mart-cont}). Procesy o prawie na pewno
ciągłych trajektoriach zastępujemy w razie potrzeby
nierozróżnialnymi wersjami o wszędzie ciągłych
trajektoriach, zmieniając je na zbiorze zerowym z $\calF_0$.

\section{Dlaczego kwadraty przyrostów są inne niż zwykła wariacja?}
Dla gładkiej funkcji $f$ na $[0,T]$ suma
$\sum_i(f(t_{i+1})-f(t_i))^2$ dąży do zera, gdy
średnica podziału dąży do zera: jest nie większa od
$\max_i|\Delta_i f|\sum_i|\Delta_i f|$, a pierwsza
wielkość maleje do zera, podczas gdy druga pozostaje
ograniczona przez wariację całkowitą. Dla ruchu Browna
typowy przyrost jest rzędu $\sqrt{\Delta t}$, zatem
jego kwadrat ma rząd $\Delta t$. Suma kwadratów
nie musi już zanikać. Właśnie ta obserwacja motywuje
rachunek Itô, który pojawi się w tomie III.

\begin{definicja}[Sumy kwadratów przyrostów]
Niech $X$ będzie rzeczywistym procesem o trajektoriach
ciągłych i $T>0$. Dla deterministycznego podziału
$\pi=\{0=t_0<t_1<\cdots<t_m=T\}$, $t\in[0,T]$,
połóżmy
\[
V_t^\pi(X)=\sum_{i=0}^{m-1}
\bigl(X_{t\wedge t_{i+1}}-X_{t\wedge t_i}\bigr)^2.
\]
Jest to ciągły proces nieujemny, choć dla ustalonej
siatki nie musi być niemalejący pomiędzy jej węzłami.
Średnicą podziału nazywamy $|\pi|=\max_i(t_{i+1}-t_i)$.
Gdy istnieje ciągły proces $A$ taki, że dla każdego
$T>0$ i każdego ciągu deterministycznych podziałów
$\pi_n$ przedziału $[0,T]$ o średnicach $|\pi_n|\to0$
mamy
$\sup_{t\le T}|V_t^{\pi_n}(X)-A_t|\to0$
według prawdopodobieństwa, proces $A$ nazywamy
\emph{wariacją kwadratową} $X$ i oznaczamy $[X]$.
Skrót \emph{ucp} oznacza właśnie zbieżność jednostajną
na przedziałach zwartych według prawdopodobieństwa.
Taka granica jest jednoznaczna z dokładnością do
nierozróżnialności: dla dwóch granic $A,A'$ na każdym
$[0,T]$ nierówność trójkąta daje
$\Prob(\sup_{t\le T}|A_t-A'_t|>\varepsilon)=0$
dla każdego $\varepsilon>0$. Wystarczy następnie
wziąć przeliczalnie wiele dodatnich całkowitych $T$
i $\varepsilon=1/n$.
Jeżeli $X$ jest adaptowany, to $V^\pi(X)$ również;
ciągłą granicę można wybrać adaptowaną, korzystając
z zupełności filtracji, jak w dowodzie poniżej.
\end{definicja}

\begin{lemat}[Znikanie wariacji kwadratowej dla trajektorii o skończonej wariacji]
\label{lem:finite-var-zero-quadratic}
Niech $f:[0,T]\to\R$ będzie funkcją ciągłą o
skończonej wariacji całkowitej $\operatorname{Var}_{[0,T]}(f)$.
Dla dowolnego ciągu podziałów o średnicach malejących
do zera
$\sup_{t\le T}V_t^{\pi_n}(f)\to0$.
\end{lemat}
\begin{proof}
Dla każdego $t\le T$ mamy
\[
V_t^\pi(f)\le
\Bigl(\max_i\sup_{u,v\in[t_i,t_{i+1}]}|f(u)-f(v)|\Bigr)
\operatorname{Var}_{[0,T]}(f).
\]
Z jednostajnej ciągłości $f$ na zwartym przedziale
pierwszy czynnik dąży do zera, jednostajnie w $t$.
Drugi czynnik nie zależy od podziału ani od $t$.
\end{proof}

\begin{twierdzenie}[Wariacja kwadratowa procesu Wienera]
\label{thm:brownian-bracket}
Niech $W$ będzie standardowym procesem Wienera,
a $T>0$. Dla dowolnych deterministycznych podziałów
$\pi_n$ o $|\pi_n|\to0$ zachodzi
$V_T^{\pi_n}(W)\to T$ w $L^2$, a nawet
\[
\sup_{0\le t\le T}|V_t^{\pi_n}(W)-t|
\longrightarrow0
\quad\text{według prawdopodobieństwa}.
\]
Ponadto dla każdego ciągu siatek spełniającego
$\sum_n|\pi_n|<\infty$ zbieżność jednostajna
zachodzi prawie na pewno. Dotyczy to w szczególności
siatek równomiernych dyadycznych.
Wynik jest zgodny z konstrukcją w rozdziale~\ref{ch:wiener}.
\end{twierdzenie}
\begin{proof}
Z niezależności przyrostów $\Delta_iW\sim
\Normal(0,h_i)$, gdzie $h_i=t_{i+1}-t_i$,
wynika $\E(\Delta_iW)^2=h_i$ i
$\Var((\Delta_iW)^2)=2h_i^2$.
Dlatego
\[
\E V_T^\pi(W)=T,\qquad
\E|V_T^\pi(W)-T|^2=2\sum_i h_i^2
\le2T|\pi|.
\]
To daje zbieżność w $L^2$ dla wszystkich siatek.
Aby uzyskać zbieżność jednostajną, w węzłach podziału
rozważmy dyskretny martyngał
$D_k=\sum_{i<k}((\Delta_iW)^2-h_i)$.
Nierówność maksymalna Dooba w $L^2$
(tom I, rozdział~\ref{I-ch:nadmartyngaly}) daje
\[
\E\max_{k\le m}|D_k|^2\le4\E|D_m|^2
\le8T|\pi|.
\]
Między węzłami różnica $V_t^\pi(W)-t$ może
zmienić się względem lewego węzła najwyżej o
$\max_i\sup_{u,v\in[t_i,t_{i+1}]}|W_u-W_v|^2+|\pi|$.
Ta wielkość dąży prawie na pewno do zera dzięki
jednostajnej ciągłości trajektorii. Stąd teza ucp.
Jeśli $\sum_n|\pi_n|<\infty$, to
$\sum_n\E\max_k|D_k|^2<\infty$
(w szczególności dla siatek dyadycznych $|\pi_n|=T2^{-n}$).
Nierówność Markowa i lemat Borela--Cantellego dają
zbieżność $\max_k|D_k|\to0$ prawie na pewno.
Oscylacje wewnątrz komórek także zanikają prawie na pewno,
więc zbieżność jest jednostajna prawie na pewno.
\end{proof}

\begin{center}
\begin{tikzpicture}[x=1.5cm,y=0.85cm,>=Latex]
\draw[->] (0,0)--(4.4,0) node[right] {$t$};
\draw[->] (0,0)--(0,3.8) node[above] {$x(t)$};
\foreach \x in {1,2,3,4}{\draw[gray,densely dotted] (\x,0)--(\x,3.3);}
\draw[blue!75!black,thick] (0,1)--(0.55,2.4)--(1,1.8)--(1.6,0.6)--(2,1.2)--(2.5,2.8)--(3,2.2)--(3.5,1.4)--(4,2.8);
\foreach \x/\y in {0/1,1/1.8,2/1.2,3/2.2,4/2.8}{\fill[blue!75!black] (\x,\y) circle (1.6pt);}
\end{tikzpicture}

\begin{tikzpicture}[x=1.5cm,y=0.85cm,>=Latex]
\draw[->] (0,0)--(4.4,0) node[right] {$t$};
\draw[->] (0,0)--(0,4.2) node[above] {$V_t^\pi(x)$};
\foreach \x in {1,2,3,4}{\draw[gray,densely dotted] (\x,0)--(\x,3.85);}
\draw[blue!75!black,thick,domain=0:0.55,samples=25] plot(\x,{(1.4*\x/0.55)^2});
\draw[blue!75!black,thick,domain=0.55:1,samples=25] plot(\x,{(1.4-0.6*(\x-0.55)/0.45)^2});
\draw[blue!75!black,thick,domain=1:1.6,samples=25] plot(\x,{0.64+(1.2*(\x-1)/0.6)^2});
\draw[blue!75!black,thick,domain=1.6:2,samples=25] plot(\x,{0.64+(1.2-0.6*(\x-1.6)/0.4)^2});
\draw[blue!75!black,thick,domain=2:2.5,samples=25] plot(\x,{1+(1.6*(\x-2)/0.5)^2});
\draw[blue!75!black,thick,domain=2.5:3,samples=25] plot(\x,{1+(1.6-0.6*(\x-2.5)/0.5)^2});
\draw[blue!75!black,thick,domain=3:3.5,samples=25] plot(\x,{2+(0.8*(\x-3)/0.5)^2});
\draw[blue!75!black,thick,domain=3.5:4,samples=30] plot(\x,{2+(-0.8+1.4*(\x-3.5)/0.5)^2});
\draw[orange!80!black,thick,dashed] (0,0)--(1,0)--(1,0.64)--(2,0.64)--(2,1)--(3,1)--(3,2)--(4,2)--(4,2.36);
\end{tikzpicture}
\par\smallskip\small Góra: przykładowa ciągła trajektoria
odcinkami liniowa i siatka $\{0,1,2,3,4\}$.
Dół: suma kwadratów z niedokończonym przyrostem (linia ciągła)
oraz suma tylko pełnych przyrostów (linia przerywana).
Pierwsza może lokalnie maleć, druga jest niemalejąca.
Nie jest to symulacja ruchu Browna.
\end{center}

\section{Dyskretny lemat izometrii: fundament konstrukcji}
Zanim skonstruujemy wariację kwadratową \emph{dowolnego}
ciągłego martyngału, potrzebujemy tylko elementarnego
faktu o skończonych sumach. Nie konstruujemy jeszcze
ogólnej całki Itô.

\begin{lemat}[Izometria dla skończonej sumy przewidywalnej]
\label{lem:discrete-isometry}
Niech $0=t_0<\cdots<t_m=T$ oraz
$M$ będzie martyngałem ciągłym i kwadratowo całkowalnym
na $[0,T]$. Jeżeli $H_i$ jest ograniczoną
$\calF_{t_i}$-mierzalną zmienną losową, to
\[
S_k:=\sum_{i=0}^{k-1}H_i(M_{t_{i+1}}-M_{t_i}),
\qquad k=0,\ldots,m,
\]
jest martyngałem dyskretnym, a
\[
\E S_m^2
=\sum_{i=0}^{m-1}\E[H_i^2(M_{t_{i+1}}-M_{t_i})^2]
=\E\sum_{i=0}^{m-1}H_i^2
\E[(M_{t_{i+1}}-M_{t_i})^2\mid\calF_{t_i}].
\]
\end{lemat}
\begin{proof}
Warunkowanie daje
$\E[H_i(M_{t_{i+1}}-M_{t_i})\mid\calF_{t_i}]=0$,
więc $S$ jest martyngałem. Po rozwinięciu kwadratu
$S_m^2$ wyrazy mieszane mają średnią zero:
jeśli $i<j$, to iloczyn
$H_i\Delta_iM H_j$ jest $\calF_{t_j}$-mierzalny,
a $\E[\Delta_jM\mid\calF_{t_j}]=0$.
Pozostają wyłącznie wyrazy diagonalne; ich druga
postać to zwykła własność warunkowej wartości oczekiwanej.
\end{proof}

\section{Konstrukcja nawiasu ciągłego martyngału bez całki Itô}
\begin{twierdzenie}[Istnienie wariacji kwadratowej]
\label{thm:general-bracket}
Niech $(\calF_t)_{t\ge0}$ będzie filtracją
spełniającą warunki zwykłe,
a $M$ ciągłym lokalnym martyngałem rzeczywistym,
$M_0=0$. Istnieje, z dokładnością do nierozróżnialności, dokładnie jeden
ciągły, adaptowany, niemalejący proces $[M]$, rozpoczynający się od zera,
taki że $M_t^2-[M]_t$ jest lokalnym martyngałem.
Dla każdego $T>0$ i dowolnej rodziny deterministycznych
podziałów $\pi_n$ przedziału $[0,T]$ o
$|\pi_n|\to0$ zachodzi
\[
\sup_{0\le t\le T}|V_t^{\pi_n}(M)-[M]_t|
\xrightarrow{\Prob}0.
\]
Dla ciągłego lokalnego martyngału o dowolnej
wartości początkowej definiujemy $[M]=[M-M_0]$;
przyrosty obu procesów są identyczne.
Wówczas $(M-M_0)^2-[M]$ jest lokalnym martyngałem.
Nie twierdzimy bez dodatkowego założenia,
że $M^2-[M]$ jest lokalnym martyngałem,
gdyż $M_0^2$ może nie być całkowalna.
\end{twierdzenie}
\begin{proof}
\emph{Krok 1: martyngał ograniczony na $[0,T]$.}
Załóżmy najpierw, że $|M_t|\le K$ dla $t\le T$.
Dla podziału $\pi$ definiujemy
\[
I_t^\pi=\sum_i M_{t_i}
(M_{t\wedge t_{i+1}}-M_{t\wedge t_i}).
\]
Dla ustalonego podziału $I^\pi$ jest ciągłym martyngałem:
każdy składnik jest zerowy przed $t_i$, a dla
$t\ge t_i$ jest martyngałem $M_{t_i}(M_{t\wedge t_{i+1}}-M_{t_i})$.
Dla $s<t_i<t$ warunkujemy najpierw względem
$\calF_{t_i}$, uzyskując zero, a następnie względem
$\calF_s$. Dla $t_i\le s\le t$ wystarcza
własność martyngałowa zatrzymanego procesu i
$\calF_s$-mierzalność współczynnika.
Ograniczenie $M$ zapewnia całkowalność.
Z dyskretnej izometrii i ortogonalności przyrostów
$\E|I_t^\pi|^2\le K^2\E M_t^2$.

Pokażemy, że $(I^\pi)$ jest ciągiem Cauchy'ego w
ucp przy malejącej średnicy podziałów. Ustalmy
$\varepsilon>0$ i dwa podziały $\pi,\rho$.
Niech $\sigma$ będzie pierwszą chwilą, w której
oscylacja $M$ od lewego końca komórki $\pi$ lub
$\rho$ osiąga poziom $\varepsilon$, oraz $T$,
jeśli nie nastąpi to wcześniej. Formalnie dla każdej komórki $[u,v]$ obu siatek
bierzemy pierwszy czas $q\in[u,v]$ taki, że
$|M_q-M_u|\ge\varepsilon$ (lub $\infty$, gdy go nie ma).
Zdarzenie trafienia do chwili $t\ge u$ wyraża się przez
$\sup_{u\le q\le v\wedge t}|M_q-M_u|\ge\varepsilon$,
które jest $\calF_t$-mierzalne dzięki ciągłości
i możliwości użycia przeliczalnego gęstego zbioru
chwil wraz z końcami przedziału.
Minimum tych czasów i $T$ jest więc czasem zatrzymania.
Z jednostajnej ciągłości trajektorii wynika
$\Prob(\sigma<T)\to0$ przy
$|\pi|\vee|\rho|\to0$.
Dla zatrzymanego martyngału $M^\sigma$ przyrosty
w każdej komórce obu siatek mają moduł co najwyżej
$\varepsilon$ (również po zatrzymaniu).
Na wspólnym rozdrobnieniu siatek różnicę
$I^\pi(M^\sigma)-I^\rho(M^\sigma)$ zapisujemy jako
skończoną sumę przewidywalną z współczynnikami
$H_j=M^\sigma_{\text{lewy }\pi}-
M^\sigma_{\text{lewy }\rho}$ o $|H_j|\le2\varepsilon$.
Z lematu~\ref{lem:discrete-isometry} oraz nierówności
maksymalnej Dooba w $L^2$ otrzymujemy
\[
\E\sup_{t\le T}|I_t^\pi(M^\sigma)-
I_t^\rho(M^\sigma)|^2
\le16\varepsilon^2\E (M^\sigma_T)^2
\le16\varepsilon^2\E M_T^2.
\]
Pierwsze oszacowanie dla supremum można otrzymać także
bezpośrednio z ciągłoczasowej nierówności Dooba;
ostatnią nierówność zapewnia próbkowanie opcjonalne
dla podmartyngału $M^2$.
Na $\{\sigma=T\}$ oba procesy zatrzymane są zgodne
z niezatrzymanym na $[0,T]$. Stąd dla $\eta>0$
\[
\Prob\left(\sup_{t\le T}|I_t^\pi-I_t^\rho|>\eta\right)
\le\Prob(\sigma<T)+16\varepsilon^2\E M_T^2/\eta^2.
\]
Najpierw wybieramy małe $\varepsilon$, potem drobne
siatki: prawa strona może być dowolnie mała.
Aby skonstruować granicę, wybieramy jeden ciąg siatek
o średnicach zbiegających do zera, np. dyadycznych.
Z własności Cauchy'ego wybieramy podciąg $(\pi_{n_k})$
tak, aby
$\Prob(\sup_{t\le T}|I^{\pi_{n_{k+1}}}_t-I^{\pi_{n_k}}_t|>2^{-k})
\le2^{-k}$.
Lemat Borela--Cantellego i zupełność przestrzeni
$C([0,T])$ z normą supremum dają jednostajną
zbieżność podciągu p.n. do ciągłej trajektorii $I$.
Poza zbiorem pełnej miary przyjmujemy zerową
trajektorię; dzięki warunkom zwykłym $I_t$ jest
$\calF_t$-mierzalna. Oszacowanie Cauchy'ego i
nierówność trójkąta pokazują, że dla każdego
innego ciągu siatek o zanikającej średnicy
granica ucp jest tym samym $I$.
Ponieważ zmienne $I_t^\pi$ są jednostajnie ograniczone
w $L^2$, zbieżność według prawdopodobieństwa implikuje
zbieżność w $L^1$ w każdej ustalonej chwili.
Przechodząc do granicy w warunkowych równościach,
stwierdzamy, że $I$ jest martyngałem.


Oszacowanie $\E|I_t^\pi|^2\le K^2\E M_t^2$
stosuje się również dla $t$ wewnątrz komórki,
ponieważ można zakończyć podział w $t$.


\emph{Krok 2: konstrukcja procesu niemalejącego.}
Z elementarnej tożsamości
$(b^2-a^2)=2a(b-a)+(b-a)^2$
na każdym odcinku podziału otrzymujemy, również
w chwilach leżących wewnątrz jego komórek,
\[
M_t^2-M_0^2=2I_t^\pi+V_t^\pi(M).
\]
Zatem $V^\pi(M)\to A:=M^2-M_0^2-2I$ w ucp.
Z ucp wybieramy podciąg podziałów, dla którego
zbieżność jest jednostajna prawie na pewno.
Proces $A$ ma wtedy trajektorie ciągłe, $A_0=0$
i $A_t\ge0$. Pokażemy, że jest niemalejący:
dla dowolnych $s<t$ proces $V^\pi$ może się obniżyć
między $s$ i $t$ tylko wskutek zmiany niedokończonego
kwadratu w komórce zawierającej $s$. Dlatego
\[
V_t^\pi(M)-V_s^\pi(M)
\ge-\max_i\sup_{u,v\in[t_i,t_{i+1}]}
|M_u-M_v|^2.
\]
Prawa strona dąży prawie na pewno do zera dzięki
jednostajnej ciągłości trajektorii. Po przejściu
do granicy $A_t\ge A_s$, najpierw dla wymiernych
$s,t$, a potem dla wszystkich dzięki ciągłości.
Ponieważ $M^2-A=2I$ jest martyngałem na $[0,T]$,
otrzymaliśmy żądany proces w przypadku ograniczonym.

\emph{Krok 3: zgodność konstrukcji na horyzontach i lokalizacja.}
Najpierw dla globalnie ograniczonego martyngału
łączymy wyniki z horyzontów $T=1,2,\ldots$.
Dla $S<T$ całkowitych porównujemy podziały $[0,T]$
zawierające $S$, których ograniczenia do $[0,S]$
mają zanikającą średnicę. Sumy na $[0,S]$ są
identyczne, więc jednoznaczność granicy ucp daje
zgodność procesów skonstruowanych na obu horyzontach.
Przeliczalność horyzontów pozwala połączyć je poza
jednym zbiorem zerowym w ciągły, adaptowany,
niemalejący proces $A$ na całej półprostej.
Proces $M^2-A$ jest martyngałem.

Dla dowolnego ciągłego lokalnego martyngału o $M_0=0$
niech $\tau_K=\inf\{t:|M_t|\ge K\}$, $K\in\mathbb N$,
$K\ge1$. Na mocy rozdziału~\ref{ch:mart-cont}
proces $M^{\tau_K}$ jest ograniczonym martyngałem,
a $\tau_K\uparrow\infty$ p.n. Oznaczmy przez $A^K$
jego skonstruowaną wariację kwadratową.
Dla każdego podziału mamy dokładną tożsamość
\[
V_t^\pi(M^{\tau_K})
=V_{t\wedge\tau_K}^\pi(M^{\tau_{K+1}}).
\]
Zbieżność ucp po prawej stronie jest zachowana
po zastąpieniu $t$ przez $t\wedge\tau_K$, ponieważ
supremum błędu nie może się zwiększyć.
Z jednoznaczności granicy ucp wynika więc
$A^K_t=A^{K+1}_{t\wedge\tau_K}$ jednocześnie
dla wszystkich $t$, poza jednym zbiorem zerowym.
Na wspólnym zbiorze pełnej miary dla wszystkich $K$
definiujemy $[M]_t$ jako stabilizującą się wartość
$A^K_t$ dla dostatecznie dużego $K$ (gdy $t\le\tau_K$).
Na dopełnieniu przyjmujemy zerową trajektorię.
Na każdym zwartym przedziale $[M]$ zgadza się z pewnym
$A^K$, więc jest ciągły, niemalejący i skończony.
Adaptacja wynika z tego, że w każdej ustalonej chwili
jest granicą $\calF_t$-mierzalnych $A^K_t$,
ze zmianą na zbiorze z $\calF_0$.
Ponadto $[M]_{t\wedge\tau_K}=A^K_t$.
W konsekwencji
\[
(M^2-[M])_{t\wedge\tau_K}
=(M^{\tau_K}_t)^2-A^K_t
\]
jest martyngałem. To dowodzi, że $M^2-[M]$
jest lokalnym martyngałem.
Dla każdego $T>0$ i $\eta>0$ mamy
\[
\Prob\left(\sup_{t\le T}|V_t^\pi(M)-[M]_t|>\eta\right)
\le \Prob(\tau_K<T)+
\Prob\left(\sup_{t\le T}|V_t^\pi(M^{\tau_K})-A^K_t|>\eta\right).
\]
Najpierw przechodzimy z $|\pi|$ do zera przy ustalonym $K$,
następnie z $K$ do nieskończoności.
Otrzymujemy zbieżność ucp dla dowolnych ciągów podziałów.

\emph{Krok 4: jednoznaczność.}
Jeśli $A$ i $A'$ spełniają tezę, to $D=A-A'$ jest
ciągłym lokalnym martyngałem (różnicą dwóch lokalnych
martyngałów po wspólnej lokalizacji)
o skończonej wariacji na przedziałach zwartych,
gdyż $A,A'$ są niemalejące i skończone.
Z twierdzenia~\ref{thm:continuous-local-finite-variation}
wynika, że $D$ jest nierozróżnialny od stałego procesu
$D_0=0$. Stąd $A=A'$ z dokładnością do nierozróżnialności.
\end{proof}

\begin{uwaga}[Co dokładnie zostało udowodnione?]
Istnienie i jednoznaczność nawiasu dla \emph{ciągłych}
lokalnych martyngałów nie wymagają w tym dowodzie
całki Itô ani pełnego twierdzenia Dooba--Meyera.
Konstrukcja używa jedynie skończonych sum przewidywalnych,
ortogonalności przyrostów, zatrzymywania i lokalizacji.
Dla procesów ze skokami trzeba osobno odróżnić
wariację kwadratową (zawierającą kwadraty skoków)
od jej przewidywalnej kompensacji; por. rozdział~\ref{ch:jump-mart}.
\end{uwaga}


\begin{twierdzenie}[Dodanie procesu o skończonej wariacji]
\label{thm:finite-var-drift-bracket}
Niech $X$ będzie procesem ciągłym mającym
wariację kwadratową w sensie ucp, a $A$
procesem ciągłym o trajektoriach lokalnie
skończonej wariacji całkowitej p.n.
Wtedy
\[
[X+A]=[X],\qquad [A]=0,
\]
a sumy mieszane $\sum_i\Delta_i^tX\Delta_i^tA$
zbiegają do zera w ucp. W szczególności wynik
dotyczy $X$ będącego ciągłym lokalnym martyngałem.
\end{twierdzenie}
\begin{proof}
Na każdym $[0,T]$ lemat
\ref{lem:finite-var-zero-quadratic} daje
$\sup_{t\le T}V_t^\pi(A)\to0$ p.n. wzdłuż
dowolnego ciągu podziałów o zanikającej średnicy.
Z nierówności Cauchy'ego--Schwarza
\[
\sup_{t\le T}\left|\sum_i\Delta_i^tX\Delta_i^tA\right|
\le
\sqrt{\sup_{t\le T}V_t^\pi(X)}
\sqrt{\sup_{t\le T}V_t^\pi(A)}.
\]
Pierwszy czynnik jest ograniczony według
prawdopodobieństwa: z ucp supremum $V^\pi(X)$
zbiega według prawdopodobieństwa do skończonego
supremum ciągłego procesu $[X]$.
Iloczyn z czynnikiem zbiegającym do zera
według prawdopodobieństwa zbiega do zera
według prawdopodobieństwa (najpierw ograniczamy
pierwszy czynnik przez dużą stałą z dużym
prawdopodobieństwem, a następnie przechodzimy
do granicy dla drugiego).
Rozwinięcie kwadratu przyrostu $X+A$ kończy dowód.
\end{proof}

\begin{przyklad}[Dodanie deterministycznego dryfu nie zmienia wariacji]
Niech $W$ będzie procesem Wienera i $b\in\R$, a
$X_t=W_t+bt$. Wtedy dla każdej siatki $\pi$ na $[0,T]$
\[
V_t^\pi(X)=V_t^\pi(W)
 +2b\sum_i\Delta_i^tW\,\Delta_i^t s
 +b^2\sum_i(\Delta_i^t s)^2,
\]
gdzie $\Delta_i^tW=W_{t\wedge t_{i+1}}-W_{t\wedge t_i}$
oraz $\Delta_i^t s=(t\wedge t_{i+1})-(t\wedge t_i)$.
Nierówność Cauchy'ego--Schwarza daje, jednostajnie w $t$,
\[
\left|\sum_i\Delta_i^tW\,\Delta_i^t s\right|
\le \sqrt{V_t^\pi(W)}\,\sqrt{T|\pi|},\qquad
\sum_i(\Delta_i^t s)^2\le T|\pi|.
\]
Z twierdzenia~\ref{thm:brownian-bracket} supremum $V_t^\pi(W)$
jest ograniczone według prawdopodobieństwa, toteż oba dodatkowe
składniki znikają w ucp. W konsekwencji $[W+bt]_t=t$.
Dryf wpływa na średnią $X_t$, ale nie na jego kwadratową
fluktuację w drobnej skali. Jest to jedno z podstawowych
rozróżnień wykorzystywanych później w rachunku It\^o.
\end{przyklad}

\begin{uwaga}[Dla ciągłego martyngału oba nawiasy się pokrywają]
Proces $[M]$ dla ciągłego lokalnego martyngału jest adaptowany
i ciągły, a zatem przewidywalny (rozdział~\ref{ch:filtracje}). Istotnie, na każdym przedziale
skończonym można go aproksymować punktowo procesami prostymi,
które na $(k2^{-n},(k+1)2^{-n}]$ przyjmują znaną już w chwili
$k2^{-n}$ wartość $[M]_{k2^{-n}}$; ich granica jest
przewidywalna jako granica funkcji mierzalnych względem
$\sigma$-ciała przewidywalnego. Różnica $(M-M_0)^2-[M]$
jest lokalnym martyngałem. W przyszłej ogólnej definicji
przewidywalnego nawiasu, opartej na kompensatorze $[M]$,
otrzymamy więc $\langle M\rangle=[M]$.
Rozbieżność tych nawiasów zaczyna mieć znaczenie przy skokach.
\end{uwaga}


\section{Kowariacja i iloczyn martyngałów}
\begin{definicja}[Kowariacja]
Dla ciągłych lokalnych martyngałów $M,N$ definiujemy
\[
[M,N]_t=\frac14\bigl([M+N]_t-[M-N]_t\bigr).
\]
Jest to ciągły proces o skończonej wariacji lokalnej.
\end{definicja}

\begin{twierdzenie}[Wzór iloczynowy dla nawiasu]
\label{thm:cross-bracket}
Niech $M,N$ będą ciągłymi lokalnymi martyngałami
z $M_0=N_0=0$. Wtedy
$M_tN_t-[M,N]_t$ jest lokalnym martyngałem oraz
$[M,N]$ jest granicą ucp sum mieszanych
$\sum_i\Delta_i^tM\Delta_i^tN$ na dowolnych
ciągach podziałów deterministycznych o średnicach
zbiegających do zera, gdzie
$\Delta_i^tM=M_{t\wedge t_{i+1}}-M_{t\wedge t_i}$
i analogicznie dla $N$.
Ponadto $[M,M]=[M]$ i
$[aM+bN,K]=a[M,K]+b[N,K]$ dla stałych $a,b$.
\end{twierdzenie}
\begin{proof}
Tożsamość
$4M_tN_t=(M_t+N_t)^2-(M_t-N_t)^2$
po odjęciu analogicznej tożsamości dla nawiasów
daje różnicę dwóch lokalnych martyngałów;
po wspólnej lokalizacji jest to martyngał.
Również dla sum po podziałach mamy
$4\sum\Delta M\Delta N=
\sum(\Delta M+\Delta N)^2-
\sum(\Delta M-\Delta N)^2$.
Zbieżność ucp obu sum kwadratów daje drugą tezę.
Dwuliniowość wynika z tej granicy, ponieważ
sumy mieszane są symetryczne i dwuliniowe dla każdego podziału,
a zatem ich granice ucp również (każda tożsamość zachodzi
z dokładnością do nierozróżnialności).
Dla $M=N$ sumy mieszane są sumami kwadratów,
co daje $[M,M]=[M]$. Wszystkie tożsamości dotyczące
przyrostów i kowariacji pozostają prawdziwe dla
dowolnych wartości początkowych, po odjęciu $M_0,N_0$.
\end{proof}

\begin{twierdzenie}[Macierz kowariacji i nierówność Kunity--Watanabego]
\label{thm:kunita-watanabe}
Niech $M=(M^{(1)},\ldots,M^{(d)})$ będzie
wektorem ciągłych lokalnych martyngałów.
Macierz $C_t=([M^{(i)},M^{(j)}]_t)_{i,j=1}^d$
jest symetryczna i jej przyrosty $C_t-C_s$
dla $s\le t$ są dodatnio półokreślone,
jednocześnie dla wszystkich $s,t$ p.n.
Dla dwóch składowych, oznaczonych $U,V$, oraz
dowolnych funkcji borelowskich $f,g$ na $[0,T]$
zachodzi trajektoryjna nierówność
\[
\int_{[0,T]} |fg|\,\mathrm d|[U,V]|
\le
\left(\int_{[0,T]} |f|^2\,\mathrm d[U]\right)^{1/2}
\left(\int_{[0,T]} |g|^2\,\mathrm d[V]\right)^{1/2},
\]
jeżeli obie całki po prawej są skończone.
Zapis $\mathrm d[U]$ oznacza miarę Lebesgue'a--Stieltjesa
trajektorii $[U]$, z zerową masą w zerze, a
$\mathrm d|[U,V]|$ oznacza wariację całkowitą
miary znakowanej indukowanej przez $[U,V]$.
W szczególności
\[
\operatorname{Var}_{[s,t]}([U,V])
\le\sqrt{([U]_t-[U]_s)([V]_t-[V]_s)}.
\]
\end{twierdzenie}
\begin{proof}
Dla $a\in\Q^d$ z dwuliniowości nawiasu mamy
$a^\top C_t a=[\sum_i a_i M^{(i)}]_t$.
Przecinając zbiory pełnej miary dla przeliczalnie
wielu $a$, otrzymujemy wszystkie te tożsamości
jednocześnie w czasie. Prawa strona jest
niemalejąca. Przez gęstość $\Q^d$ i ciągłość
formy kwadratowej w $a$ przyrosty $C_t-C_s$
są dodatnio półokreślone także dla dowolnego
$a\in\R^d$.

Ustalmy trajektorię z tego zbioru pełnej miary,
a na $[0,T]$ oznaczmy przez $\mu,\nu,\gamma$
miary Lebesgue'a--Stieltjesa odpowiednio
$[U],[V],[U,V]$. Dwie pierwsze są nieujemne
i skończone, trzecia znakowana i skończonej
wariacji. Dla dowolnych $a,b\in\R$ miara
$a^2\mu+2ab\gamma+b^2\nu$ jest nieujemna,
gdyż odpowiada niemalejącej funkcji $[aU+bV]$.
Stąd dla każdego zbioru borelowskiego $E$
macierz z elementami $\mu(E),\gamma(E),\nu(E)$
jest dodatnio półokreślona, a więc
$|\gamma(E)|^2\le\mu(E)\nu(E)$.
W szczególności $\gamma\ll\kappa=\mu+\nu$.
Z twierdzenia Radona--Nikodyma oznaczmy
$p=\mathrm d\mu/\mathrm d\kappa$,
$q=\mathrm d\nu/\mathrm d\kappa$,
$h=\mathrm d\gamma/\mathrm d\kappa$.
Mamy $p,q\ge0$, $p+q=1$ $\kappa$-p.w.
Dla wymiernych $a,b$ nieujemność powyższej
miary daje $a^2p+2abh+b^2q\ge0$ $\kappa$-p.w.
Po przecięciu przeliczalnej rodziny zbiorów
pełnej miary i przejściu do rzeczywistych $a,b$
dostajemy $|h|^2\le pq$.
Wariacja całkowita $\gamma$ ma gęstość $|h|$
względem $\kappa$ (wynika to z rozkładu na
dodatnią i ujemną część gęstości).
Zatem
\[
\int|fg|\,\mathrm d|\gamma|
\le\int |f|\sqrt p\,|g|\sqrt q\,\mathrm d\kappa
\le
\left(\int|f|^2\,\mathrm d\mu\right)^{1/2}
\left(\int|g|^2\,\mathrm d\nu\right)^{1/2}
\]
przez nierówność Cauchy'ego--Schwarza.
Dla $f=g=\chi_{(s,t]}$ otrzymujemy oszacowanie
wariacji na $[s,t]$, ponieważ miary nie mają atomów
z uwagi na ciągłość funkcji.
Argument jest trajektoryjny, więc obowiązuje
także dla funkcji zależnych od $\omega$,
o ile dla danej trajektorii są borelowskie.
\end{proof}

\begin{wniosek}[Kowariacja ruchu Browna]
Dla standardowego ruchu Browna
$B=(B^{(1)},\ldots,B^{(d)})$ zachodzi
$[B^{(i)},B^{(j)}]_t=\delta_{ij}t$.
Jeżeli $X_t=AB_t$ dla deterministycznej macierzy $A$,
to $[X^{(i)},X^{(j)}]_t=(AA^\top)_{ij}t$.
\end{wniosek}
\begin{proof}
Dla $i=j$ stosujemy twierdzenie
\ref{thm:brownian-bracket}. Dla $i\ne j$ niezależność
współrzędnych daje, dla ustalonego podziału,
\[
\E\left|\sum_k\Delta_kB^{(i)}\Delta_kB^{(j)}\right|^2
=\sum_k (t_{k+1}-t_k)^2\le T|\pi|\to0.
\]
Granica w $L^2$ jest zero, a z poprzedniego
twierdzenia musi być równa kowariacji. Zależność
dla $AB$ wynika z dwuliniowości.
\end{proof}


\section{Warunkowa wariancja przyrostu i elementarna izometria}
\begin{twierdzenie}[Wariancja warunkowa przez nawias]
\label{thm:conditional-variance-bracket}
Niech $M$ będzie ciągłym martyngałem na $[0,T]$
z $M_T\in L^2$ oraz $M_0=0$.
Wtedy $[M]_T\in L^1$, a $M^2-[M]$ jest
martyngałem na $[0,T]$.
Dla każdych $0\le s\le t\le T$ zachodzi
\[
\E[(M_t-M_s)^2\mid\calF_s]
=\E[[M]_t-[M]_s\mid\calF_s].
\]
W szczególności $\E M_t^2=\E\bigl([M]_t\bigr)$.
\end{twierdzenie}
\begin{proof}
Dla $\tau_K=\inf\{t\ge0:|M_t|\ge K\}$
(z konwencją $\inf\varnothing=\infty$ na $[0,T]$)
proces $M^{\tau_K}$ jest ograniczonym martyngałem.
Z twierdzenia~\ref{thm:general-bracket} i własności
zatrzymywania nawiasu (twierdzenie~\ref{thm:stopped-bracket})
wynika, że
$M_{t\wedge\tau_K}^2-[M]_{t\wedge\tau_K}$
jest martyngałem na $[0,T]$, rozpoczynającym się
od zera. W szczególności
$\E[M]_{t\wedge\tau_K}=\E M_{t\wedge\tau_K}^2$.
Z nierówności Dooba (twierdzenie~\ref{thm:doob-cont})
$\E\sup_{u\le T}|M_u|^2\le4\E M_T^2<\infty$.
Kwadraty zatrzymanych wartości zbiegają więc w $L^1$
do $M_t^2$ przez zbieżność zmajoryzowaną.
Ponieważ $[M]$ jest niemalejący i $\tau_K$ rosną,
zbieżność monotoniczna daje
\[
\E[M]_t=\lim_K\E[M]_{t\wedge\tau_K}=\E M_t^2<\infty.
\]
W szczególności $[M]_{t\wedge\tau_K}\to[M]_t$
w $L^1$. Przechodząc do granicy w warunkowych
równościach dla zatrzymanych martyngałów,
uzyskujemy własność martyngałową $M^2-[M]$.
Zatem
$\E[M_t^2-M_s^2\mid\calF_s]
=\E[[M]_t-[M]_s\mid\calF_s]$.
Ponieważ $M_t-M_s$ jest warunkowo wycentrowany,
$\E[2M_s(M_t-M_s)\mid\calF_s]=0$.
Rozwinięcie $(M_t-M_s)^2=M_t^2-M_s^2-2M_s(M_t-M_s)$
daje tezę. Biorąc $s=0$, otrzymujemy równość średnich.
\end{proof}

\begin{wniosek}[Warunkowa kowariancja przyrostów]
\label{cor:conditional-cross-bracket}
Dla ciągłych martyngałów $M,N$ kwadratowo
całkowalnych na $[0,T]$ z zerowymi wartościami
początkowymi proces $MN-[M,N]$ jest martyngałem i
\[
\E[(M_t-M_s)(N_t-N_s)\mid\calF_s]
=\E[[M,N]_t-[M,N]_s\mid\calF_s]
\]
dla $0\le s\le t\le T$. W szczególności
$\E M_tN_t=\E[M,N]_t$.
\end{wniosek}
\begin{proof}
Stosujemy twierdzenie
\ref{thm:conditional-variance-bracket}
do $M+N$ i $M-N$ oraz odejmujemy
odpowiednie równości, dzieląc przez cztery.
Całkowalność wynika również z nierówności
Cauchy'ego--Schwarza dla iloczynu i z
nierówności~\ref{thm:kunita-watanabe} dla wariacji
kowariacji, gdyż $\E[M]_T,\E[N]_T<\infty$.
\end{proof}

\begin{wniosek}[Izometria elementarnych strategii]
\label{cor:elementary-bracket-isometry}
Przy założeniach poprzedniego twierdzenia niech
$0=t_0<\cdots<t_n=T$ i $H_i$ będą ograniczonymi
$\calF_{t_i}$-mierzalnymi zmiennymi. Wtedy
\[
\E\left(\sum_{i=0}^{n-1}H_i(M_{t_{i+1}}-M_{t_i})\right)^2
=\E\sum_{i=0}^{n-1}H_i^2
([M]_{t_{i+1}}-[M]_{t_i}).
\]
\end{wniosek}
\begin{proof}
Stosujemy dyskretną izometrię z lematu
\ref{lem:discrete-isometry}. Dla każdej komórki
zastępujemy warunkowy drugi moment przyrostu przez
warunkowy przyrost $[M]$, zgodnie z poprzednim
twierdzeniem. Ponieważ $H_i^2$ jest
$\calF_{t_i}$-mierzalne, można je wyciągnąć przed
warunkowanie, a następnie wziąć wartość oczekiwaną.
\end{proof}

\begin{wniosek}[Izometria mieszana elementarnych strategii]
\label{cor:elementary-cross-isometry}
Przy założeniach wniosku~\ref{cor:conditional-cross-bracket}
niech $H_i,K_i$ będą ograniczonymi
$\calF_{t_i}$-mierzalnymi zmiennymi.
Wówczas
\[
\E\left[
\left(\sum_i H_i\Delta_i M\right)
\left(\sum_i K_i\Delta_i N\right)
\right]
=\E\sum_i H_iK_i
([M,N]_{t_{i+1}}-[M,N]_{t_i}).
\]
\end{wniosek}
\begin{proof}
Po rozwinięciu iloczynu wyrazy dla różnych
indeksów mają średnią zero, warunkując względem
początku późniejszego przyrostu.
Wyrazy diagonalne przekształcamy przez
wniosek~\ref{cor:conditional-cross-bracket}
i wyciągnięcie $H_iK_i$ przed warunkowanie.
Całkowalność wszystkich wyrazów zapewniają
ograniczenie współczynników, kwadratowa
całkowalność martyngałów oraz oszacowanie
wariacji z twierdzenia~\ref{thm:kunita-watanabe}.
\end{proof}

\begin{uwaga}[Gdzie w tej formule ukrywa się całka Itô?]
Lewa strona opisuje skończoną sumę przewidywalnych
zakładów, których stawka $H_i$ jest znana przed
przyrostem $M_{t_{i+1}}-M_{t_i}$. Prawa strona
mierzy ich kwadratowy koszt, ważony przyrostami
wariacji kwadratowej. W tomie III rozszerzymy tę
równość z procesów prostych na odpowiednie klasy
procesów przewidywalnych przez domknięcie podprzestrzeni w odpowiedniej przestrzeni
Hilberta; właśnie wtedy pojawi się całka stochastyczna.
\end{uwaga}

\begin{uwaga}[Dlaczego zwykła wariacja i nawias opisują inne zjawiska?]
Dla funkcji ciągłej $f$ o skończonej wariacji całkowitej
suma $\sum_i(f(t_{i+1})-f(t_i))^2$ zanika przy
zagęszczaniu podziału, niezależnie od tego, jak duża
jest całkowita zmiana $f$. Dla ruchu Browna te same
sumy dążą do długości przedziału. Nawias jest więc
miarą fluktuacji drugiego rzędu, a nie sumą bezwzględnych
przyrostów. W szczególności ciągły lokalny martyngał
może mieć skończony nawias na $[0,T]$, mimo że jego
trajektorie mają nieskończoną wariację całkowitą.
\end{uwaga}

\section{Wariacja kwadratowa a czas zatrzymania}
\begin{twierdzenie}[Wariacja kwadratowa procesu zatrzymanego]
\label{thm:stopped-bracket}
Jeśli $M$ jest ciągłym lokalnym martyngałem i
$\tau$ czasem zatrzymania, to
$[M^\tau]_t=[M]_{t\wedge\tau}$ dla $t\ge0$.
W szczególności dla procesu Wienera
$[W^\tau]_t=t\wedge\tau$.
\end{twierdzenie}
\begin{proof}
Dla każdego deterministycznego podziału $\pi$ i każdej
chwili $t$ zachodzi dokładna równość
$V_t^\pi(M^\tau)=V_{t\wedge\tau}^\pi(M)$.
Wynika to bezpośrednio z definicji: po chwili $\tau$
proces zatrzymany jest stały, a w komórce zawierającej
$\tau$ ostatni przyrost kończy się w $\tau$.
Z twierdzenia~\ref{thm:local-mart-criteria} proces
$M^\tau$ pozostaje lokalnym martyngałem.
Z twierdzenia~\ref{thm:general-bracket} obie strony
mają granice ucp, odpowiednio $[M^\tau]_t$ oraz
$[M]_{t\wedge\tau}$: dla prawej strony supremum
błędu po zatrzymaniu nie przekracza supremum
błędu na całym $[0,T]$. Jednoznaczność granicy według
prawdopodobieństwa daje równość procesów.
Przecięcie zdarzeń pełnej miary dla chwil wymiernych
i ciągłość nawiasów zapewniają jednoczesność dla wszystkich $t$.
\end{proof}

\begin{zadanie}[Brownowski martyngał i czas wyjścia z przedziału]
Niech $W$ będzie procesem Wienera,
$\tau=\inf\{t\ge0:|W_t|=a\}$ dla $a>0$,
a $M_t=W_{t\wedge\tau}$.
Wyznaczyć $[M]_t$ oraz $\E[M_t^2]$.

\emph{Rozwiązanie.} Z twierdzenia o zatrzymaniu nawiasu
$[M]_t=t\wedge\tau$. Ponieważ $|M_t|\le a$,
$M$ jest ograniczonym martyngałem, a
$M_t^2-[M]_t$ martyngałem o średniej zero przy $t=0$.
Dlatego $\E M_t^2=\E(t\wedge\tau)$.
Gdy $t\to\infty$, lewa strona dąży do $a^2$,
a prawa do $\E\tau$ dzięki zbieżności monotonicznej.
Otrzymujemy ponownie $\E\tau=a^2$, bez używania
całki Itô.
\end{zadanie}


% ================================================================
% ROZDZIAŁ 14. SKOKI, KOMPENSATORY, PODSUMOWANIE TOMU
% ================================================================
